\documentclass[12pt]{article}
\usepackage[utf8]{inputenc}
\usepackage[T1]{fontenc}
\usepackage[brazil]{babel}
\usepackage[a4paper,margin=2.5cm]{geometry}
\usepackage{enumitem}
\setlist[itemize]{noitemsep,topsep=4pt}
\title{Sprint 3 - S3\_05\\Busca e Ordenação}
\author{Matheus Moreli Durão Silva\\Matrícula 12621CBS217\\Cibersegurança - Bacharelado - Noturno}
\date{2026}
\begin{document}
\maketitle

\section*{Entregável}
Implementações e comparação de busca e ordenação. Este documento registra o objetivo do card, a solução desenvolvida e a forma de conferência. Os arquivos de código e as evidências citadas permanecem na mesma pasta deste documento.

\section*{O que foi estudado e aplicado}
\begin{itemize}
\item Busca binária exige dados ordenados
\item Bubble, Selection e Insertion Sort
\item Análise de complexidade
\item Vantagens e limitações
\end{itemize}

\section*{Como conferir}
O avaliador pode abrir os arquivos desta pasta, executar os exemplos quando houver código e comparar a saída com a explicação apresentada. Nos cards ligados ao Inventário de Ativos, a solução foi relacionada ao Trabalho 2 sem modificar requisitos já atendidos do CRUD.

\section*{Comparação reproduzível}
\texttt{comparar.py} conta comparações em listas de 20, 100 e 300 itens com ordem crescente, inversa e embaralhada (semente 42). Para 100 itens, Bubble/Selection/Insertion fazem respectivamente 99/4950/99 comparações na entrada ordenada, 4950/4950/4950 na inversa e 4940/4950/2597 na aleatória. A busca binária pelo último item usa 7 comparações após ordenação. Para 300 itens, os valores aleatórios são 44772/44850/24468, e a busca usa 9.
\section*{Complexidade e escolha}
Bubble com parada antecipada e Insertion são $O(n)$ no melhor caso e $O(n^2)$ no pior; Selection é $O(n^2)$ nos dois. Bubble e Insertion são estáveis, Selection por troca pode não ser. Busca binária custa $O(\log n)$, mas exige a lista previamente ordenada; para uma consulta em lista desordenada, busca linear pode ser melhor. O CRUD principal consulta IDs por dicionário. A tabela completa, os limites e o script reprodutível estão em \texttt{relatorio\_comparativo.md}.
\section*{Conclusão}
As contagens resultam de código executado sobre entradas fixas, com as saídas conferidas contra \texttt{sorted}.
\end{document}
